% SOURCE_AUTHORITY: sga5_fr_workpass.tex, lines 8797--8856 (LF-normalized unit).
% SOURCE_SHA256: 791F4EFFC5E02832D5D77ED03518C8156D6F07E4C8238B03545DB93D883FBB28
\subsubsection*{2.2. Imagens diretas superiores}

Seja
\[
f:X\longrightarrow Y
\]
um morfismo separado de tipo finito de pré-esquemas localmente noetherianos. Recordemos o resultado seguinte de cohomologia étale (SGA 4 XVII; para o caso próprio, cf. SGA 4 XIV 1.1):

\begin{statement}{Lema 2.2.1}
Sejam \(A\) um anel comutativo anulado por um inteiro \(n\) não nulo e \(M\) um \((A,\Ab(X))\)-módulo construtível. Para todo inteiro \(i\), o feixe
\[
\R^i_!f(M)
\]
é um \((A,\Ab(Y))\)-módulo construtível.
\end{statement}

O lema seguinte deduz-se aplicando (V 5.3.1) após restringir aos abertos noetherianos de \(Y\):

\begin{statement}{Lema 2.2.2}
Para todo inteiro \(i\) e todo \((AR,\mathbb Z_\ell)\)-feixe construtível \(F\) sobre \(X\), o sistema projetivo
\[
\R^i_!f(F)
\]
é um \((AR,\mathbb Z_\ell)\)-feixe construtível.
\end{statement}

Dito isto, o \(\delta\)-funtor exato
\[
(\R^i_!f)_{i\in\mathbb Z}:\Ab(X)\longrightarrow \Ab(Y)
\]
se prolonga de maneira evidente a um \(\delta\)-funtor exato, denotado da mesma forma,
\[
\Hom(\mathbb N^\circ,\Ab(X))\longrightarrow \Hom(\mathbb N^\circ,\Ab(Y)),
\]
que fornece por restrição um \(\delta\)-funtor exato
\[
U(X)\longrightarrow U(Y).
\tag{1}
\]
Por outro lado, os funtores \(\R^i_!f\) transformam um sistema projetivo localmente \(AR\)-nulo em um sistema projetivo localmente \(AR\)-nulo, de modo que, por passagem ao quociente, o \(\delta\)-funtor (1) fornece um \(\delta\)-funtor exato% CORRECCIÓN DE FUENTE n.º 60: notación «R^i_!f» (! bajo R) armonizada en seis apariciones de §2.1--2.2 (L8808--8849); el libro, pp. 274--275, pone R^i_! f (! bajo R) en todas partes, al igual que el resto del .tex (L8861+, todo XV); estas seis contenían «R^if_!» (! bajo f), en desacuerdo con el libro y con la convención dominante. Verificación a 200 dpi de pp. 274 y 275. 2026-07-03.
\[
(\R^i_!f)_{i\in\mathbb Z}:(AR,\mathbb Z_\ell)\text{-}\fc(X)
\longrightarrow (AR,\mathbb Z_\ell)\text{-}\fc(Y).
\]
Usando agora a equivalência de categorias (1.5.5), este último \(\delta\)-funtor é interpretado como um \(\delta\)-funtor exato
\[
\mathbb Z_\ell\text{-}\fc(X)\longrightarrow \mathbb Z_\ell\text{-}\fc(Y).
\]

\begin{statement}{Definição 2.2.3}
O funtor
\[
\R^i_!f:\mathbb Z_\ell\text{-}\fc(X)\longrightarrow \mathbb Z_\ell\text{-}\fc(Y)
\]
chama-se \(i\)-ésimo funtor de imagem direta com suportes próprios de \(f\).
\end{statement}

\begin{statement}{Observação}
Supondo demonstrado o análogo do lema (2.2.1) para as imagens diretas superiores ordinárias quando \(X\) e \(Y\) são excelentes (o que está feito em característica zero (SGA 4 XIX 5.1)), definem-se do mesmo modo, neste caso, imagens diretas superiores em cohomologia \(\ell\)-ádica. Poder-se-ia desenvolver igualmente, neste quadro, o resultado do número seguinte.
\end{statement}
